\subsection{Sparse operator algebra, symmetrization, gauge correction, and density normalization}
\label{subsec:postprocessing}

Sparse operator manipulations are implemented directly on aligned block data.
Given two sparse operators $A$ and $B$ with reverse-edge alignment already
established, \texttt{Mandala} evaluates their trace contraction as
\begin{equation}
\Tr(AB)
=
\sum_{(\bm{L},i,j)}
\mathrm{tr}\!\left(A_{ij}^{\bm{L}} B_{ji}^{-\bm{L}}\right),
\end{equation}
without assembling dense global matrices. This operator algebra is reused for
energy and electron-count evaluation, loss construction, and diagnostic
analysis.

For matrix targets expected to be symmetric or Hermitian in the chosen basis,
\texttt{Mandala} supports explicit sparse symmetrization after block reconstruction. In
the real-valued convention used by the implementation, this takes the
form
\begin{equation}
X_{ij}^{\bm{L}}
\leftarrow
\frac{1}{2}
\left(
X_{ij}^{\bm{L}} + \left(X_{ji}^{-\bm{L}}\right)^{\top}
\right).
\end{equation}
This operation is carried out using precomputed reverse-edge permutations, so
that symmetry is enforced consistently across periodic images and self-edges.

Hamiltonian gauge freedom is handled operationally through the overlap-shift
correction described in Section~\ref{subsubsec:gauge_invariance}. In practical
evaluation and analysis workflows, \texttt{Mandala} computes
\begin{equation}
\widetilde{H}^{\mathrm{pred}}
=
H^{\mathrm{pred}} - \mu_H S,
\qquad
\mu_H
=
\frac{
\left\langle H^{\mathrm{pred}} - H^{\mathrm{ref}}, S \right\rangle_F
}{
\left\langle S, S \right\rangle_F
},
\end{equation}
so that purely gauge-like Hamiltonian offsets are removed before reporting
gauge-corrected operator errors.

The framework also supports density normalization to the correct number of
electrons. If a predicted density $\widehat{D}$ and overlap $\widehat{S}$
produce
$\widehat{N}_e = \Tr(\widehat{D}\widehat{S})$, then a normalized density is
obtained by the scalar rescaling
\begin{equation}
\widehat{D}_{\mathrm{norm}}
=
\beta \widehat{D},
\qquad
\beta = \frac{N_e^{\mathrm{target}}}{\Tr(\widehat{D}\widehat{S})}.
\end{equation}
This provides a simple and effective way to enforce electron-count consistency.

Together, sparse trace evaluation, symmetry restoration, gauge correction, and
density normalization define a compact postprocessing layer that turns raw
network outputs into physically interpretable operators ready for observable
evaluation and Hamiltonian-based analyses such as band-structure calculation.
